\chapter*{\centering Abstract}
\onehalfspacing
Low-rank adaptation (LoRA) fine-tunes a pretrained model by training small low-rank matrices next to its frozen
weights. This report measures what LoRA costs and saves against full fine-tuning for two vision transformers,
ViT-B/16 and Swin-B, on CIFAR-100, in the same configurations on both models and over five seeds. We ask four
questions: is LoRA as accurate as full fine-tuning, what does it save during training, where should the adapters be
placed, and what does it change when the model is deployed.

With the learning rate of each method chosen on validation loss, LoRA at rank 8 is more accurate than full
fine-tuning on Swin-B (by 0.4--0.5 percentage points for each of three choices of adapted layers). On ViT-B it is
within about 0.2 points of full fine-tuning, with no evidence that it is better. Its saving in training memory depends on which layers are adapted,
in the same way on both models: about a third for the attention projections, a seventh for the MLP layers and about 6\%
for both, at batch size 64. The saving grows as the batch shrinks, to about 70\% at batch size 1. Training epochs
are 12--26\% shorter. Within LoRA, the training memory is set mainly by the first block that has an adapter, because
the backward pass must reach it; adapting only the second half of the blocks needs the least memory but loses
1.2--1.6 points of accuracy. For deployment, merging the adapters leaves the predictions unchanged and gives the
memory and speed of an ordinary model, while adapters kept separate slow down single-image predictions by 8--55\% on
the GPU. A closing study shows that peak training memory can be predicted from the shape of the model before
measuring it (within 1.5\% on ViT-L), whereas a simple rule for the accuracy of a placement, written down before
the runs, did not hold.
